\chapter{Reality as Constraint}

\section{What Makes a Trajectory Real?}

Chapter 37 closed by noting that not every trajectory Chapter 13's proposal machinery could generate is one this book's architecture would actually permit into a world's history. C, the constraint manifold Chapter 18 built for engineering reasons, does the work of picking out which trajectories are admissible and which are not. This chapter asks what that work actually amounts to, once it is examined philosophically rather than only operationally: what is admissibility doing, when it does the deciding between a trajectory that becomes part of a world and one that does not?

\section{Possibility, Actuality, and the Write Cycle}

The traditional vocabulary for this question is modal: possibility and actuality, what could be the case against what actually is the case. This book's write cycle, examined again with this vocabulary in mind, turns out to map onto that distinction with unusual precision. A proposal, \(\Delta\psi\), is a possibility in close to the philosopher's sense: a candidate way things might go, generated but not yet evaluated against anything. Projection, \(\Pi_C\), is the evaluation, checking that candidate against whatever the world's own constraints actually permit. Commit is actualization itself, the specific, well-defined operation by which a mere possibility becomes part of what a world's history actually contains. What is worth noticing is how unusual it is for a philosophical account of modality to come with an actual mechanism for this last step. Most treatments of possibility and actuality leave the transition from one to the other either unexplained or grounded in something this book has no occasion to invoke, a metaphysically primitive act of actualization, an observer collapsing a wave function, the best of all possible worlds being selected by some cosmic optimization. This architecture supplies, instead, a precisely specified operation, checkable and, in the engineering chapters, actually implementable, and it is worth taking seriously that this is not nothing, whatever else remains open about how far the analogy to physical modality extends.

\section{Refusal as Genuine Impossibility}

This gives modal vocabulary a further, more precise distinction this book's own machinery already computes without needing new philosophical apparatus. A proposal that \(\Pi_C\) can correct into some nearby admissible state, section 14.3's ordinary case, was possible in a qualified sense: not actual exactly as proposed, but actual in some modified form close enough to preserve its intent. A proposal that triggers Chapter 15.4's refusal, where no admissible state honors what was proposed closely enough to count as a correction rather than a substitution, is impossible in a considerably stronger sense: not merely a possibility that failed to occur, but a candidate the world's own constraints rule out entirely. This book's architecture does not treat possibility and impossibility as a binary applied from outside, by an observer's judgment about what seems conceivable. It computes the distinction directly, as a consequence of whether an admissible nearby state to a given proposal actually exists.

\section{Is C Arbitrary?}

It is worth confronting directly a tension the last several chapters have left unaddressed. Chapter 35 argued that world-building, in fiction, is constraint design: an author genuinely chooses what C contains, magic, faster-than-light travel, talking animals, entirely at their own discretion. If C can simply be chosen, does treating admissibility as a serious account of reality collapse into something arbitrary, reality being whatever anyone happens to stipulate?

The honest answer requires distinguishing what varies across this book's applications from what does not. C's source genuinely differs by domain: authorial choice in fiction, per Chapter 35; empirical discovery in physical simulation, per Chapter 32, where C is not chosen but measured, constrained by what actual experiments reveal about actual conservation laws; something closer to negotiated, emergent convention in the institutions of Chapter 36, where a society's laws are neither arbitrarily stipulated by one author nor simply read off from physical measurement, but arrived at through processes this book has not attempted to formalize. For the case that matters most to a metaphysical claim, the actual universe's own C, this book offers no account of why it contains what it contains rather than something else. This is not a gap unique to this framework. It is one of the oldest open questions in the philosophy of physics, why the laws are these laws and not others, and this book does not resolve it. What this chapter can honestly claim is narrower and still worth having: that once reality is understood as admissibility plus persistence, the question of where a given domain's C actually comes from becomes a well-posed, meaningful question with a real answer to be sought, rather than a confusion to be dissolved. Fiction's C is chosen because fiction's whole point permits choosing it. Physics's C, whatever it ultimately turns out to be grounded in, is emphatically not chosen in that sense, and this book's framework has enough structure to at least state that difference precisely, even where it cannot yet explain it.

\section{Existence Without Substance}

This closes the loop back to Chapter 37 tightly enough to be worth stating as a single combined claim. Existence, under this framework, is not an intrinsic property something has by virtue of being made of the right kind of stuff. It is a structural, relational property: a configuration exists, in the full sense this book has been developing, when it satisfies whatever C actually governs its domain and reproduces itself as a genuine fixed point or conserved orbit within that constraint. Substance metaphysics asked what a thing is made of. This framework asks instead whether a candidate trajectory is admissible and whether, once admitted, it sustains itself. Both questions can be asked of the same wall. Only the second one, this book has argued across thirty-seven chapters, was ever doing real explanatory work.

\section{Looking Ahead}

This chapter has treated reality, admissibility, and existence at a level of generality broad enough to apply to a courtyard, a robot's gait, or a civilization's institutions alike. It has not yet asked the more specific and more exposed question of whether this vocabulary can do the work of physics itself, actual physical law, rather than serving only as a useful borrowing of physics-flavored mathematics for domains that are not, strictly, physical at all. Chapter 39 asks that question directly.
